% template for poster sections 
% use thie file to create individual sections for a beamer poster 
% then put the pieces together using the poster-template file 

%%%%%%%%%%%%%%%%%%%% 
% top matter 
%%%%%%%%%%%%%%%%%%%% 

%\documentclass[leqno,presentation]{beamer} 
\documentclass[leqno,handout]{beamer}

%\usepackage[orientation=portrait, size=a3, scale=1]{beamerposter} 
\DeclareGraphicsExtensions{.eps, .jpg, .png, .tif} 

\usepackage{calc} 
\usepackage{times} 
\usepackage{type1cm} 
\usepackage[latin1]{inputenc} 
\usepackage{xcolor, colortbl} 
\usepackage{amsmath, amssymb, latexsym} 
\usepackage[english]{babel} 
\usetheme{Darmstadt} 
\usepackage{hyperref} 
\usepackage{gensymb}


% change to location of graphics files, comment out if graphics files 
% are in same directory as tex file 

\graphicspath{{assets/}} 

%%%%%%%%%%%%%%%%%%%%%%%%%%%%%%%%%%%%%%%% 
% end top matter, document begins 
%%%%%%%%%%%%%%%%%%%%%%%%%%%%%%%%%%%%%%%% 

\newcommand{\highlight}[1]{ 
\addtolength{\fboxrule}{.2ex} 
\begin{block}{} 
\begin{quote}#1 
\end{quote} 
\end{block} 
} 

\title{Symmetric Fixed Points of the Ricci Flow} 
	\author{\small{\\[-2.5ex] Faculty Mentor: Eric Chen \\
 	Project Leader: Anthony D'Arienzo \\
	IML Scholars: Noah Kang}\\}

\institute{\\[-3ex] University of Illinois at Urbana-Champaign 
  
           \\[3ex] Illinois Mathematics Lab \\ 
           \\[2ex] \includegraphics[width = 0.3\textwidth]{uiuc_logo-with-name.png} \\
           [2ex]
                 
                   Midterm Presentation \\ 
                   October 15, 2026\\ [2ex]} 
\date{} 
\begin{document} 
\frame{\titlepage} 

%%%%%%%%%%%%%%%%%%%% 

\begin{frame} 
\frametitle{Introduction} 
\begin{itemize}
    \item \textbf{Ricci flow} is a process that smooths out curved shapes over time, but it can sometimes cause them to collapse into "singularities".
    \item To fix these singularities, we use special shapes called \textbf{Expanding Ricci Solitons} that grow in a perfectly uniform, self-similar way.
    \item They are defined by a metric $g$ and a potential function $f$ satisfying: $\text{Ric} + \nabla^2f + g = 0$.
    \item We are studying three specific families of these fixed shapes: $R^3\times S^1$, $R^2\times S^2$, and $R^n$ (Bryant solitons).
\end{itemize}
\end{frame} 

%%%%%%%%%%%%%%%%%%%% 

\begin{frame} 
\frametitle{Project Goals} 
\begin{itemize}
    \item \textbf{Goal:} Explore and prove how these fixed shapes grow and behave and study the properties as we go to spatial infinity (i.e. what the shapes/spaces look like as our distance from a fixed reference points becomes very large).
    \item \textbf{Question 1:} What shape do these objects approach as they grow infinitely large? (Their "asymptotic cone")
    \item \textbf{Question 2:} How do the starting conditions dictate the final, infinite-size shape?
    \item \textbf{Question 3:} For Bryant solitons, how does the starting state exactly determine the final shape?
\end{itemize}
\end{frame} 

%%%%%%%%%%%%%%%%%%%% 

\begin{frame} 
\frametitle{Governing ODEs} 
\begin{itemize}
    \item By assuming the shapes have high symmetry, we can reduce the complex math into a system of Ordinary Differential Equations (ODEs):
    \begin{align*}
        a'' &= -2\frac{a'b'}{b} + a'f' + a \\
        b'' &= \frac{1 - (b')^2}{b} - \frac{a'b'}{a} + b'f' + b \\
        f'' &= \frac{a''}{a} + 2\frac{b''}{b} - 1
    \end{align*}
    \item \textbf{Challenge:} These equations require dividing by variables that start at zero ($b=0$). This creates a mathematical "division-by-zero" error right at the starting line ($r=0$).
\end{itemize}
\end{frame} 

%%%%%%%%%%%%%%%%%%%% 

\begin{frame} 
\frametitle{Computer Assisted Proof} 
\begin{itemize}
  \item Goal: We want to deduce from our numerics the genuine existence of expanders asymptotic to cones with $R>0$.
  \item $r=\epsilon \to r=0.01$: Found max error of Taylor Series expansion from 11th term and beyond.
  \item $r=0.01 \to r=20$: Used a Julia script to bound errors for $a,a',b,b',f,f'$ from $r=0.01$ to $r=20$.
  \item $r=20\to r=\infty$: Define Lyapunov Ratios to verify that the space expands smoothly without singularity for all $r \in [0,\infty)$.
  
\end{itemize}
\end{frame} 

%%%%%%%%%%%%%%%%%%%% 

\begin{frame}
\frametitle{Results: Mapping the Asymptotic Cones (3D View)}
\begin{figure}
    \centering
    \includegraphics[width=10\textwidth,height=0.65\textheight,keepaspectratio]{f2_map_asymptotic_manifold.png}
    \caption{Graphed $b_0$ on the z-axis with $a'_{\infty}$ and $b'_{\infty}$ on the x and y axes and $f''(0)$ color-scaled.}
\end{figure}
\end{frame}

%%%%%%%%%%%%%%%%%%%%

\begin{frame}
\frametitle{Results: Achievable Cone Limits (2D Projection)}
\begin{figure}
  \centering
  \includegraphics[width=0.85\textwidth,height=0.65\textheight,keepaspectratio]{f2_projection.png}
  \caption{Projected previous image onto $a'_{\infty}$ and $b'_{\infty}$ as x and y axes.}
\end{figure}
\end{frame}

%%%%%%%%%%%%%%%%%%%%

\begin{frame}
\frametitle{Results: Bryant Soliton Thresholds}
\begin{figure}
    \centering
    \includegraphics[width=0.8\textwidth,height=0.55\textheight,keepaspectratio]{asymptotic_cone_bryan_soliton.png}
    \caption{For Case 2.4, we plotted $\lim_{r\to \infty} \sqrt{h(r)}$ on the y axis and $\mu_1=h'(0)$ on the x axis.}
\end{figure}
\end{frame}

%%%%%%%%%%%%%%%%%%%% 

\begin{frame} 
\frametitle{Next Steps} 
\begin{itemize}
    \item Estimate the relative entropies of expanders, defined by Deruelle-Schulze, asymptotic to the same cone", where relative entropy has been defined.
\end{itemize}
\end{frame} 

%%%%%%%%%%%%%%%%%%%% 

\end{document}